\section{平展态射}\label{sec:ag-3-3}
% 原PDF p50起.
回顾，所考虑的概形均为 Noether 概形，态射则为有限型态射或局部态射.
\begin{numpara}[凝聚层支集的回顾]\label{item:ag-3-3-1}
设 $\mathcal M$ 是 $X$ 上的凝聚层. 称使纤维 $\mathcal M_x$ 非零的点 $x\in X$ 所成的集合为 $\mathcal M$ 的\newterm{ag-support}{支集}{support}，记为 $\operatorname{Supp}(\mathcal M)$. 它是 $X$ 的闭集. 当 $X=\operatorname{Spec}(A)$，且 $\mathcal M$ 由模 $M$ 给出时，$\mathcal M$ 的支集就是零化子理想 $\operatorname{Ann}(M)$ 定义的闭集 $V(\operatorname{Ann}(M))$.
\end{numpara}
\begin{definition}\label{item:ag-3-3-2}
若 $\Omega_{Y/X}$ 在 $y$ 处的纤维 $\Omega_{Y/X,y}$ 为零，则称态射 $f:Y\longrightarrow X$ 在点 $y\in Y$ 处\newterm{ag-unramified-morphism}{非分歧}{unramified morphism}. 若它在每一点处都非分歧，即 $\Omega_{Y/X}=0$，则称它非分歧.
\end{definition}
\begin{remark}\label{item:ag-3-3-3}
\begin{enumerate}[label=(\roman*)]
\item 可见，一个态射的非分歧点集是开集. 此开集可能为空，域的不可分扩张便给出了这样的例子.
\item 由命题~\ref{item:ag-3-2-9}，非分歧态射的复合仍非分歧. 由推论~\ref{item:ag-3-2-5}，非分歧态射的基变换仍非分歧.
\item 浸入（开浸入或闭浸入）都是非分歧的.
\item $f$ 在 $y$ 处非分歧当且仅当对角浸入 $\Delta_{Y/X}$ 在 $y$ 的某个邻域上是开浸入. 事实上，若闭浸入由 Noether 局部环 $A$ 的理想 $I$ 给出，且余法模 $I/I^2$ 为零，则此闭浸入是同构（由 Nakayama 引理，$I$ 为零）. 更整体地说，一个态射非分歧当且仅当相应的对角浸入是开浸入. 特别地，非分歧态射总是在命题~\ref{item:ag-3-2-9}意义下微分光滑的.\unskip\nobreak\hfill\ensuremath{\diamond}
\end{enumerate}
\renewcommand{\envenddiamond}{}
\end{remark}
\begin{proposition}\label{item:ag-3-3-4}
有限型态射 $f:Y\longrightarrow X$ 在点 $y$ 处非分歧，当且仅当令 $x:=f(y)$ 后，有 $\mathfrak m_x\mathcal O_{Y,y}=\mathfrak m_y$，且 $\kappa(y)$ 是 $\kappa(x)$ 的有限可分扩张.
\end{proposition}
\begin{proof}
因 $\Omega$ 的构造与局部化相容，有
\begin{align*}
\Omega_{Y/X,y}=\Omega_{\mathcal O_{Y,y}/\mathcal O_{X,x}}.
\end{align*}
再将推论~\ref{item:ag-3-2-5}应用于由 $\operatorname{Spec}(\kappa(x))\hookrightarrow\operatorname{Spec}(\mathcal O_{X,x})$ 给出的基变换，得到
\begin{align*}
\Omega_{\mathcal O_{Y,y}/\mathcal O_{X,x}}\otimes_{\mathcal O_{Y,y}}\kappa(y)=\Omega_{(\mathcal O_{Y,y}/\mathfrak m_x\mathcal O_{Y,y})/\kappa(x)}\otimes\kappa(y).
\end{align*}
由 Nakayama 引理，$\mathcal O_{Y,y}$-模 $\Omega_{\mathcal O_{Y,y}/\mathcal O_{X,x}}$ 为零，当且仅当 $\kappa(y)$-向量空间 $\Omega_{\mathcal O_{Y,y}/\mathcal O_{X,x}}\otimes_{\mathcal O_{Y,y}}\kappa(y)$ 为零，故由上式，当且仅当 $\Omega_{(\mathcal O_{Y,y}/\mathfrak m_x\mathcal O_{Y,y})/\kappa(x)}$ 为零.
因此，只须证明：若 $A$ 是局部 $k$-代数，则 $\Omega_{A/k}=0$ 当且仅当 $A$ 是 $k$ 的可分扩张. 方向 $\Leftarrow$ 已经证明.
现假设 $\Omega_{A/k}=0$，先证明 $K:=A/\mathfrak m$ 是 $k$ 的有限可分扩张. 满射 $A/\mathfrak m\otimes_A\Omega_{A/k}\twoheadrightarrow\Omega_{(A/\mathfrak m)/k}$ 说明 $\Omega_{K/k}=0$. 将 $K$ 写成某个 $k$ 上函数域 $k(X_i)_{i\in I}$ 的可分扩张 $K'$ 的纯不可分扩张. 正合列
\begin{align*}
\Omega_{K'/k}\otimes K\longrightarrow\Omega_{K/k}\longrightarrow\Omega_{K/K'}\longrightarrow0
\end{align*}
迫使 $\Omega_{K/K'}=0$，所以 $K'=K$，因为已假设 $K'$ 是纯不可分的. 因此 $K$ 在 $k(X_i)_{i\in I}$ 上可分，故态射 $\operatorname{Spec}(K)\longrightarrow\operatorname{Spec}(k(X_i))$ 非分歧，从而微分光滑. 它显然也平坦，故由命题~\ref{item:ag-3-2-9}得到同构 $\Omega_{k(X_i)/k}\otimes K\xrightarrow\sim\Omega_{K/k}$. 将已算出的仿射空间余切模局部化，得到 $\dim_K(\Omega_{K/K})=\operatorname{card}(I)$，即超越基的基数. 由于要求此维数为零，故 $K$ 在 $k$ 上代数可分. 又因假设中所有态射都是有限型的，此扩张还是有限的.
还须证明 $A=K$，即 $A$ 的根基 $\mathfrak m$ 为零. 设 $\bar k$ 是 $k$ 的代数闭包，只须证明 $A_{\bar k}:=A\otimes_k\bar k$ 的根基为零，或者说，它的每个局部化（其数目有限）都是域. 有 $\Omega_{A_{\bar k}/\bar k}=0$，因而对 $A_{\bar k}$ 的每个局部化 $A'$，有 $\Omega_{A'/\bar k}$. $A'$ 的剩余域是 $\bar k$ 的有限可分扩张，故等于 $\bar k$. 例~\ref{item:ag-3-2-2}说明 $\mathfrak m'/\mathfrak m'^2\simeq A'/\mathfrak m'\otimes_{A'}\Omega_{A'/\bar k}$，所以 $\mathfrak m'/\mathfrak m'^2=0$，由 Nakayama 引理得到 $\mathfrak m'=0$.
\end{proof}
\begin{remark}\label{item:ag-3-3-5}
上述命题表明，这种几何意义下的分歧概念与算术中的分歧概念相容. 设 $F$ 是 $\mathbb Q$ 的有限扩张，$\mathcal O_F$ 为其整数环. 态射 $Y=\operatorname{Spec}(\mathcal O_F)\longrightarrow X=\operatorname{Spec}(\mathbb Z)$ 的分歧点集，就是 $\mathcal O_F$ 中位于某个按算术意义``在 $F$ 中分歧''的素数 $p$ 上方的素理想全体. 算术中使用的``差理想''，正是 $\Omega_{Y/X}$ 的零化子 $\mathfrak d_{Y/X}$.\unskip\nobreak\hfill\ensuremath{\diamond}
\renewcommand{\envenddiamond}{}
\end{remark}
\begin{numpara}[拟有限态射]\label{item:ag-3-3-6}
上述命题说明，非分歧态射是\newterm{ag-quasi-finite-morphism}{拟有限态射}{quasi-finite morphism}，这里的意思是，对 $Y$ 的每一点，$\mathcal O_{Y,y}/\mathfrak m_{f(y)}\mathcal O_{Y,y}$ 在 $\kappa(f(y))$ 上有限. 从几何角度看，这等价于要求 $y$ 是其所在纤维的孤立点（习题）. 此条件还等价于完备化 $\widehat{\mathcal O}_{Y,y}$ 在 $\widehat{\mathcal O}_{X,f(y)}$ 上有限（习题）. 而且，在此情形下，$\mathfrak m_y$-进拓扑与 $\mathfrak m_{f(y)}\mathcal O_{Y,y}$-进拓扑相同，因而典范态射
\begin{align*}
\widehat{\mathcal O}_{X,f(y)}\otimes_{\mathcal O_{X,f(y)}}\mathcal O_{Y,y}\longrightarrow\widehat{\mathcal O}_{Y,y}
\end{align*}
是同构. 由此可见，前一命题中的 $\mathcal O_{Y,y}$ 可换成它的完备化. 换言之，一点处的非分歧性可从完备局部环读出.
对有限型态射，Grothendieck 形式的 Zariski 主定理给出如下整体刻画：$f$ 拟有限当且仅当存在分解 $Y\xrightarrow iY'\xrightarrow{f'}X$，其中 $i$ 是开浸入，$f'$ 是有限态射. 关于这些概念，一份很好的参考文献是 Raynaud 的 \emph{Anneaux locaux henséliens} 第 3 章.
\end{numpara}
\begin{definition}\label{item:ag-3-3-7}
若态射 $f:Y\longrightarrow X$ 在点 $y\in Y$ 处非分歧且平坦，则称其在 $y$ 处\newterm{ag-etale-morphism}{平展}{étale morphism}. 若它在每一点处都平展，则称其平展.
\end{definition}
\begin{remark}\label{item:ag-3-3-8}
\begin{enumerate}[label=(\roman*)]
\item $f$ 的平展点集是 $Y$ 的开集.
\item 平展态射的复合及基变换仍平展.
\item 开浸入是平展态射.
\item 一个态射平展当且仅当它平坦，且相应的对角浸入是开浸入（同时也是闭浸入）.
\item 平展态射 $f:Y\longrightarrow X$ 平坦，且与所有非分歧态射一样微分光滑，所以由正合列~\ref{item:ag-3-2-9}可知，若 $X$、$Y$ 是 $S$-概形，则典范态射 $f^*\Omega_{X/S}\longrightarrow\Omega_{Y/S}$ 是同构.\unskip\nobreak\hfill\ensuremath{\diamond}
\end{enumerate}
\renewcommand{\envenddiamond}{}
\end{remark}
\begin{example}\label{item:ag-3-3-9}
设 $A$ 是环，$P\in A[T]$ 是首一多项式，且理想 $(P,P')$ 是单位理想 $A$. 则态射 $Y=\operatorname{Spec}(A[T]/(P))\longrightarrow X=\operatorname{Spec}(A)$ 平展. 事实上，正合列~\ref{item:ag-3-2-12}给出
\begin{align*}
\mathcal N_{Y/\mathbb A_X^1}\simeq\mathcal O_Y\cdot\overline P\longrightarrow\mathcal O_Y\cdot dT\longrightarrow\Omega_{Y/X}\longrightarrow0,
\end{align*}
其中 $\overline P$ 是 $P$ 在 $(P)/(P^2)$ 中的像. 第一个态射将 $\overline P$ 映到 $P'(T)dT$. 由假设，$P'$ 在 $\mathcal O_Y$ 中可逆，所以第一个态射满射，从而 $\Omega_{Y/X}=0$. 此外，因 $P$ 首一，$A[T]/(P)$ 在 $A$ 上自由，以 $i<\deg(P)$ 时 $T^i$ 的像为一组基.\unskip\nobreak\hfill\ensuremath{\diamond}
\renewcommand{\envenddiamond}{}
\end{example}
下面将看到，此例很好地代表了平展态射的局部结构. 我们需要如下平坦性判别法.
\begin{proposition}\label{item:ag-3-3-10}
设 $A$ 是 Noether 局部环，$I$ 是真理想，$M$ 是 $A$-模. 假设 $M$ 还是某个 Noether $A$-代数 $B$ 上的有限模. 则下列性质等价：
\begin{enumerate}[label=\roman*)]
\item $M$ 在 $A$ 上平坦.
\item $M/IM$ 在 $A/I$ 上平坦，且典范态射 $(M/IM)\otimes_A\operatorname{gr}^{I}(A)\longrightarrow\operatorname{gr}^{I}(M)$ 是同构.
\item 对每个 $n>0$，$M/I^nM$ 在 $A/I^n$ 上平坦.
\end{enumerate}
\end{proposition}
\begin{proof}
\textup{i)$\Rightarrow$ii).} 对每个 $n$，态射 $I^n\otimes_A M\longrightarrow I^nM$ 是单射，模 $I$ 约化后便得到结合分次模的同构. 此外，标量扩张保持平坦性，故 $M/IM$ 在 $A/I$ 上平坦.
\textup{ii)$\Rightarrow$iii).} 固定 $n$. 因为在 $A/I^n$ 中工作，可假设 $I^n=0$. 要证在假设 ii) 下，$M$ 在 $A$ 上平坦，即对每个 $A$-模 $N$，有 $\operatorname{Tor}_1^A(M,N)=0$. 用逐步分解的方法，只须对满足 $IN=0$ 的 $N$ 证明这一点；此时有正合列
\begin{align*}
\operatorname{Tor}_1^A(M,A/I)\otimes_{A/I}N\longrightarrow\operatorname{Tor}_1^A(M,N)\longrightarrow\operatorname{Tor}_1^{A/I}(M/IM,N)\longrightarrow0.
\end{align*}
由假设，右端项为零，所以只须证明 $\operatorname{Tor}_1^A(M,A/I)=0$. 这等价于映射 $I\otimes M\longrightarrow IM$ 为单射. 对 $k$ 作递降归纳，证明 $I^k\otimes M\longrightarrow I^kM$ 为单射. 当 $k=n$ 时显然成立. 从 $k+1$ 推到 $k$ 时，使用图
\begin{align*}
\begin{tikzcd}[ampersand replacement=\&]
I^{k+1}\otimes_A M\arrow[r]\arrow[d,hook] \& I^k\otimes_A M\arrow[r]\arrow[d] \& \operatorname{gr}^k A\otimes_A M\arrow[d,"\sim"] \\
I^{k+1}M\arrow[r] \& I^kM\arrow[r] \& \operatorname{gr}^k M.
\end{tikzcd}
\end{align*}
\textup{iii)$\Rightarrow$i).} 正是在这里要用到 $M$ 在中间的 Noether 代数上有限这一假设. 它保证 $M\otimes_A N$ 型的模的 $I$-进滤过是分离的. 现设 $N\hookrightarrow N'$ 是 $A$-模的单态射. 要证明 $\alpha:M\otimes_A N\longrightarrow M\otimes_A N'$ 为单射. 在如下交换图中，
\begin{align*}
\begin{tikzcd}[ampersand replacement=\&,column sep=large]
M\otimes_A N\arrow[r]\arrow[d,"\alpha"'] \& M\otimes_A(N/N\cap I^nN')\arrow[d] \\
M\otimes_A N'\arrow[r] \& M\otimes_A(N'/I^nN')
\end{tikzcd}
\end{align*}
第二条竖直箭头可等同于单射 $N/(N\cap I^nN')\hookrightarrow N'/I^nI'$ 在 $A/I^n$ 上与 $M/I^nM$ 的张量积，故由假设，它是单射. 因而 $\alpha$ 的核包含在 $M\otimes_A(N\cap I^n)$ 的（像）中. 由 Artin--Rees 引理，对每个 $m\in\mathbb N$，存在 $n$ 使得 $I^nN'\cap N\subset I^mN$. 于是 $\alpha$ 的核对每个 $m\in\mathbb N$ 都包含在 $I^m(M\otimes_A N)$ 中，因此为零.
\end{proof}
\begin{proposition}\label{item:ag-3-3-11}
设 $f:Y\longrightarrow X$ 是态射，$y\in Y$ 的像为 $x\in X$. 下列性质等价：
\begin{enumerate}[label=\roman*)]
\item $f$ 在 $y$ 处平展.
\item 典范态射 $\kappa(y)\otimes_{\kappa(x)}\operatorname{gr}^{\mathfrak m_x}(\mathcal O_{X,x})\longrightarrow\operatorname{gr}^{\mathfrak m_y}(\mathcal O_{Y,y})$ 是同构，且扩张 $\kappa(y)/\kappa(x)$ 有限可分.
\item 存在首一多项式 $Q\in\widehat{\mathcal O}_{X,x}[T]$，使得 $Q$ 与 $Q'$ 生成 $\mathcal O_{X,x}[T]$ 的单位理想，且 $\widehat{\mathcal O}_{Y,y}\simeq\widehat{\mathcal O}_{X,x}[T]/(Q)$.
\item 存在首一多项式 $P\in\mathcal O_{X,x}[T]$，使得 $\mathcal O_{Y,y}$ 同构于 $\mathcal O_{X,x}[T]/(P)$ 的一个局部化，且若以 $u$ 表示 $T$ 在 $\mathcal O_{Y,y}$ 中的像，则元素 $P'(u)$ 在 $\mathcal O_{Y,y}$ 中可逆.
\end{enumerate}
\end{proposition}
\begin{proof}
\textup{i)$\Leftrightarrow$ii).} ii) 中的典范态射是余法层对图
\begin{align*}
\begin{tikzcd}[ampersand replacement=\&]
\operatorname{Spec}(\kappa(y))\arrow[r,hook]\arrow[d] \& Y\arrow[d,"f"] \\
\operatorname{Spec}(\kappa(x))\arrow[r,hook] \& X
\end{tikzcd}
\end{align*}
的函子性的特例. 假设 $f$ 在 $y$ 处平展. 非分歧性判别法~\ref{item:ag-3-3-4}说明此图是笛卡尔图. $f$ 的平坦性保证所需的同构性质，参见~\ref{item:ag-3-1-4}. 反过来，假设分次模之间的态射是同构. 则前面的平坦性判别法保证 $f$ 在 $y$ 处平坦. 此外，特别得到箭头 $\mathcal O_{Y,y}\otimes\mathfrak m_x\longrightarrow\mathfrak m_y$ 满射，即 $\mathfrak m_x\mathcal O_{Y,y}=\mathfrak m_y$.
\textup{iii)$\Leftrightarrow$i).} 假设 iii) 成立，与前面的例子一样，$\widehat{\mathcal O}_{Y,y}$ 在 $\widehat{\mathcal O}_{X,x}$ 上平坦且非分歧. 利用命题~\ref{item:ag-3-3-4}的刻画，我们已注意到，这时 $\mathcal O_{Y,y}$ 在 $\mathcal O_{X,x}$ 上也非分歧. 上述平坦性判别法又表明，$\mathcal O_{Y,y}$ 在 $\mathcal O_{X,x}$ 上平坦. 因此 $f$ 在 $y$ 处平展. 反过来，假设 i) 成立，则 $\widehat{\mathcal O}_{Y,y}$ 在 $\widehat{\mathcal O}_{X,x}$ 上平坦且非分歧. 关键在于，$\widehat{\mathcal O}_{Y,y}$ 在 $\widehat{\mathcal O}_{X,x}$ 上还是有限的，其次数为 $\dim_{\kappa(x)}(\kappa(y))$. 取 $\widehat{\mathcal O}_{Y,y}$ 中元素 $u$，使其像 $\bar u$ 生成 $\kappa(x)$ 的有限可分扩张 $\kappa(y)$. 由 Nakayama 引理，$u$ 生成的子代数 $\widehat{\mathcal O}_{X,x}[u]$ 等于 $\widehat{\mathcal O}_{Y,y}$. 更确切地说，$u^i$（$0\leqslant i<\dim_{\kappa(x)}(\kappa(y))$）构成 $\widehat{\mathcal O}_{Y,y}$ 在 $\widehat{\mathcal O}_{X,x}$ 上的一组基. 在此基下表示 $u^{\dim_{\kappa(x)}(\kappa(y))}$，便得到可分多项式 $P\in\widehat{\mathcal O}_{X,x}[T]$（即 $(P,P')$ 是 $\mathcal O_{X,x}[T]$ 的单位理想），使 $T\mapsto u$ 诱导同构 $\widehat{\mathcal O}_{X,x}[T]/(P)\xrightarrow\sim\widehat{\mathcal O}_{Y,y}$.
\textup{iv)$\Rightarrow$iii).} 由于 $\widehat{\mathcal O}_{X,x}$ 是 Hensel 环（即满足 Hensel 引理，参见 Raynaud，\emph{Anneaux locaux henséliens} 第 1 章），$\widehat{\mathcal O}_{X,x}/(P)$ 是若干局部代数的积，而 $\widehat{\mathcal O}_{Y,y}$ 可等同于 $\widehat{\mathcal O}_{X,x}/(P)$ 的一个直因子，其形式为 $\widehat{\mathcal O}_{X,x}/(Q)$，其中 $Q\in\widehat{\mathcal O}_{X,x}[T]$ 整除 $P$. 此时 $Q'(u)$ 在 $\widehat{\mathcal O}_{Y,y}$ 中可逆.
\textup{i)$\Rightarrow$iv).} （证明梗概）由假设，即 $\mathcal O_{Y,y}$ 在 $\mathcal O_{X,x}$ 上平坦且非分歧. 与 i)$\Rightarrow$iii) 的本质区别在于，$\mathcal O_{Y,y}$ 在 $\mathcal O_{X,x}$ 上未必有限（即使 $f$ 有限也如此）. 使用前面回顾的 Grothendieck 形式的 Zariski 主定理，可知 $\mathcal O_{Y,y}$ 是某个有限 $\mathcal O_{X,x}$-代数 $B$ 在极大理想 $\mathfrak n$ 处的局部化. 可以找到 $B$ 中的一个元素 $u$，使其属于 $B$ 的其余极大理想，且其像生成 $\kappa(x)$ 的可分扩张 $B/\mathfrak n$. 然后证明，$B$ 的子代数 $\mathcal O_{X,x}[u]$ 在 $\mathfrak n\cap\mathcal O_{X,x}[u]$ 处的局部化，与 $B$ 在 $\mathfrak n$ 处的局部化同构；再证明存在 $P$，满足 $P(u)=0$、$P'(u)\in\mathfrak n$，且 $T\mapsto u$ 给出同构 $\mathcal O_{X,x}[T]/(P)\xrightarrow\sim\mathcal O_{X,x}[u]\subset B$. （见 SGA 1，第 I 讲，引理 7.9.）
\end{proof}
\begin{remark}\label{item:ag-3-3-12}
有几个特殊情形值得强调.
\begin{enumerate}[label=\roman*)]
\item 若 $\kappa(x)=\kappa(y)$，则 $f$ 平展当且仅当它诱导同构 $\widehat{\mathcal O}_{X,x}\xrightarrow\sim\widehat{\mathcal O}_{Y,y}$.
\item 查看证明 iv)$\Rightarrow$iii) 可知，所用到的完备化的唯一性质是它为 Hensel 环. 因此，iii) 中的 $\widehat{\mathcal O}_{X,x}$ 可以换成一个局部 Hensel 子 $\mathcal O_{X,x}$-代数 $B$（并将 $\widehat{\mathcal O}_{Y,y}$ 换成 $B\otimes_{\mathcal O_{X,x}}\mathcal O_{Y,y}$）. 可以证明，存在最小的这样的子代数 $B=\mathcal O_{X,x}^h$，称为 $\mathcal O_{X,x}$ 的\newterm{ag-henselization}{Hensel 化}{Henselization}. 此时有 $\mathcal O_{Y,y}^h=\mathcal O_{X,x}^h\otimes_{\mathcal O_{X,x}}\mathcal O_{Y,y}$. 参见 Raynaud，\emph{Anneaux locaux henséliens} 第 8 章.
\item 设 $Y$、$X$ 是复簇. 可以证明，$\widehat{\mathcal O}_{X,x}$、$\widehat{\mathcal O}_{Y,y}$ 中由全纯函数芽组成的子代数 $\mathcal O_{X,x}^{\mathrm{hol}}$、$\mathcal O_{Y,y}^{\mathrm{hol}}$ 是 Hensel 环（这是隐函数定理的推论，参见 Raynaud，\emph{Anneaux locaux henséliens} 第 7 章）. 在此情形下，$f$ 在 $y$ 处平展当且仅当它诱导同构 $\mathcal O_{X,x}^{\mathrm{hol}}\xrightarrow\sim\mathcal O_{Y,y}^{\mathrm{hol}}$.\unskip\nobreak\hfill\ensuremath{\diamond}
\end{enumerate}
\renewcommand{\envenddiamond}{}
\end{remark}
\begin{exercise*}
设 $f:\mathbb A_{\mathbb C}^1\longrightarrow\mathbb A_{\mathbb C}^1$ 由 $z\mapsto z^2$ 定义. 验证 $f$ 平坦，且分歧点集为 $\{0\}$. 取 $y=1$（从而 $x=1$）. 具体地验证 $\mathcal O_{Y,1}$ 在 $\mathcal O_{X,1}$ 上不有限，且嵌入 $\mathcal O_{X,1}\otimes_{\mathcal O_X}\mathcal O_Y\hookrightarrow\mathcal O_{Y,1}\times\mathcal O_{Y,-1}$ 不满射. 再验证，嵌入 $\widehat{\mathcal O}_{X,1}\otimes_{\mathcal O_X}\mathcal O_Y\hookrightarrow\widehat{\mathcal O}_{Y,1}\times\widehat{\mathcal O}_{Y,-1}$ 是双射，且 $\widehat{\mathcal O}_{X,1}\xrightarrow\sim\widehat{\mathcal O}_{Y,1}$.
\end{exercise*}
\begin{corollary}[性质的保持]\label{item:ag-3-3-13}
设态射 $f:Y\longrightarrow X$ 在点 $y\in Y$ 处平展，$y$ 的像为 $x\in X$. 则：
\begin{enumerate}[label=\roman*)]
\item $\mathcal O_{X,x}$ 与 $\mathcal O_{Y,y}$ 的维数相同，深度也相同.
\item $\mathcal O_{X,x}$ 是正则环、Cohen--Macaulay 环、正规环、既约环，当且仅当 $\mathcal O_{Y,y}$ 分别具有相同性质.
\end{enumerate}
\end{corollary}
\begin{proof}
\textup{(i)} 维数相等由前一命题的 ii) 借助 Hilbert 多项式得到.
\textup{(ii)} 正则性的保持由 (i) 得到. 关于 Cohen--Macaulay 环的断言由深度相等得到. 对正规性，使用 Serre 判别法：Noether 局部环正规当且仅当
\begin{enumerate}[label=\roman*)]
\item 对每个高度为 $1$ 的素理想 $\mathfrak p$，局部化 $A_{\mathfrak p}$ 正则.
\item 对每个高度 $\geqslant2$ 的素理想 $\mathfrak p$，$A_{\mathfrak p}$ 的深度 $\geqslant2$.
\end{enumerate}
再注意，态射 $\operatorname{Spec}(B)\longrightarrow\operatorname{Spec}(A)$ 保持理想的高度，且此态射的局部化都是平展的.
\end{proof}
对概形 $X$ 的一点 $x$，定义其\newterm{ag-tangent-space}{切空间}{tangent space}为余切空间的对偶 $(\mathfrak m_x/\mathfrak m_x^2)^*$. 它是 $\kappa(x)$ 上的向量空间. 每个将点 $y$ 映到点 $x$ 的态射 $f:Y\longrightarrow X$ 都诱导切映射
\begin{align*}
T_y(f):(\mathfrak m_y/\mathfrak m_y^2)^*\longrightarrow(\mathfrak m_x/\mathfrak m_x^2)^*\otimes\kappa(y).
\end{align*}
它是余切映射的对偶，而余切映射正是前一命题 (i) 中 $\mathfrak m$-进结合分次模之间态射的次数 $1$ 部分. 因此，从此 (i) 可抽出如下性质：
\begin{corollary}\label{item:ag-3-3-14}
对平展态射，上述切映射 $T_y(f)$ 是同构.
\end{corollary}
\begin{remark}\label{item:ag-3-3-15}
正如这一结果以及前面注中的 iii) 所提示的，平展态射是拓扑中``局部同构''在代数几何中的类似物. 对 $\mathbb C$ 上的簇，这不只是类比. 事实上，平展态射的局部形式，以及可用于参数化多项式方程解的隐函数定理，说明两个 $\mathbb C$ 上的簇之间的平展态射对于解析拓扑是局部同构（但对于 Zariski 拓扑不是）.
这一类比对代数几何的发展与应用取得了丰硕成果. 举三个例子.
\begin{enumerate}[label=\roman*)]
\item 正如局部平凡覆盖可用于定义拓扑空间的基本群，平展覆盖（见下文）可用于定义概形的基本群.
\item 已知，拓扑空间的奇异上同调可从局部平凡覆盖构成的``意象''上的常值层上同调恢复. 同样，平展态射组成一个``景''，在其上可以研究层上同调. 平展上同调具有出色的性质，尤其对域 $k$ 上的簇如此. 当 $k=\mathbb C$ 时，它恢复奇异上同调. 当 $k$ 是有限域时，它可以用于计算簇的 $k$-点个数. 对其他域 $k$，例如 $\mathbb Q$，平展上同调空间带有 $k$ 的绝对 Galois 群作用，因此具有重要的算术意义.
\item 与局部平凡覆盖一样，平展态射在形变下是刚性的. 对有限型态射而言，这甚至是一项刻画性质.\unskip\nobreak\hfill\ensuremath{\diamond}
\end{enumerate}
\renewcommand{\envenddiamond}{}
\end{remark}
鉴于这些应用的重要性，我们再作一些说明，但不详述. 定义概形基本群的方法如下（证明见 SGA 1，第 V 讲）.
\begin{definition}\label{item:ag-3-3-16}
称既平展又有限的态射为\newterm{ag-etale-cover}{平展覆盖}{étale cover}.
\end{definition}
设 $X$ 是概形，$\mathcal C_X$ 为 $X$ 的平展覆盖 $f:Y\longrightarrow X$ 所成的范畴. 有如下结果：
\begin{theorem}\label{item:ag-3-3-17}
范畴 $\mathcal C_X$ 是``\newterm{ag-galois-category}{Galois 范畴}{Galois category}''，即具有下列性质：
\begin{enumerate}[label=\roman*)]
\item 存在有限归纳极限与有限投射极限.
\item 每个态射 $f$ 都可分解为复合 $i\circ p$，其中 $p$ 为严格满态射，$i$ 是到某个直因子的同构.
\end{enumerate}
\end{theorem}
若列
\begin{align*}
\begin{tikzcd}[ampersand replacement=\&]
X\times_{X'}X\arrow[r,shift left]\arrow[r,shift right] \& X\arrow[r,"p"] \& X'
\end{tikzcd}
\end{align*}
正合，则称满态射 $p:X\longrightarrow X'$ 为\newterm{ag-strict-epimorphism}{严格满态射}{strict epimorphism}.
现设 $\bar x:\operatorname{Spec}(\Omega)\hookrightarrow X$ 为 $X$ 的一个``\newterm{ag-geometric-point}{几何点}{geometric point}''，即 $\Omega$ 是代数闭域. 定义``在 $\bar x$ 处的纤维''函子
\begin{align*}
\begin{aligned}
F_{\bar x}:\mathcal C_X&\longrightarrow\mathcal{Ens},\\
(Y\xrightarrow fX)&\longmapsto f^{-1}(\bar x).
\end{aligned}
\end{align*}
这是一个取值于有限集合的反变函子.
\begin{theorem}\label{item:ag-3-3-18}
函子 $F_{\bar x}$ 是 Galois 范畴 $\mathcal C_X$ 上的一个``\newterm{ag-fiber-functor}{纤维函子}{fiber functor}''，即具有下列性质：
\begin{enumerate}[label=\roman*)]
\item 与有限投射极限及有限归纳极限交换.
\item 若 $F_{\bar x}(f)$ 为双射，则 $f$ 为同构.
\end{enumerate}
\end{theorem}
现在应用适用于任何带纤维函子 $F$ 的 Galois 范畴 $\mathcal C$ 的抽象方法. 设 $\operatorname{Aut}(F)$ 为纤维函子 $F$ 的自同构群. 在其上赋予由核 $\operatorname{Aut}(F)\longrightarrow\operatorname{Aut}(F(Y))$（其中 $Y$ 为 $\mathcal C$ 的对象）生成的拓扑. 特别地，开子群的指数均有限. 由构造，函子 $F$ 诱导函子
\begin{align*}
\widetilde F:\mathcal C\longrightarrow\operatorname{Aut}(F)\text{-}\mathcal{Ens},
\end{align*}
其中右端是带有 $\operatorname{Aut}(F)$ 连续作用的有限集合所成的范畴. 可以证明：
\begin{proposition}\label{item:ag-3-3-19}
函子 $\widetilde F$ 是范畴等价.
\end{proposition}
在我们关心的情形中，将 $\operatorname{Aut}(F_{\bar x})$ 记为 $\pi_1(X,\bar x)$，称为带点概形 $(X,\bar x)$ 的\newterm{ag-fundamental-group}{基本群}{fundamental group}. 与拓扑情形不同，这里得到的是 投射有限拓扑群，而非离散群. 可以证明，若 $\bar y$ 是 $X$ 的另一个几何点，则 $F_{\bar y}$ 与 $F_{\bar x}$ 同构.\footnote{可以将同构 $F_{\bar x}\xrightarrow\sim F_{\bar y}$ 所成的集合，视为拓扑空间两点 $x$、$y$ 之间路径的同伦类集合的类似物.}因此，基本群不依赖基点，差别仅在于一个非唯一的同构. 当 $\pi_1(X)=0$ 时，称 $X$ \newterm{ag-simply-connected}{单连通}{simply connected}. 这等价于 $X$ 的每个平展覆盖都同构于若干个 $X$ 的副本之和.
\begin{example}\label{item:ag-3-3-20}
\begin{enumerate}[label=(\roman*)]
\item 当 $X$ 是 $\mathbb C$ 上的代数簇时，可以证明，上述代数 $\pi_1$ 可等同于拓扑 $\pi_1$ 的 投射有限完备化.
\item 若 $X=\operatorname{Spec}(k)$，其中 $k$ 是域，$\bar x=\operatorname{Spec}(\bar k)$，其中 $\bar k$ 是 $k$ 的代数闭包，则 $\pi_1(X,\bar x)=\operatorname{Gal}(\bar k/k)$.
\item 算术中可以证明，每个数域在 $\mathbb Q$ 上都至少在一个素数处分歧. 这说明 $\operatorname{Spec}(\mathbb Z)$ 没有非平凡平展覆盖，即 $\operatorname{Spec}(\mathbb Z)$ 单连通.\unskip\nobreak\hfill\ensuremath{\diamond}
\end{enumerate}
\renewcommand{\envenddiamond}{}
\end{example}
\begin{exercise}\label{item:ag-3-3-21}
验证态射 $\mathbb A_{\mathbb F_p}^1\longrightarrow\mathbb A_{\mathbb F_p}^1$，$x\mapsto x^p-x$ 是一个非平凡平展覆盖. 特别地，$\mathbb F_p$ 上的仿射直线不是单连通的！
\end{exercise}